\chapter{The Preserved Source Audit}\label{ch04:audit}

The single most consequential decision in the CAS-Ohm 1.0 candidate is a
negative one: the decision \emph{not} to harmonize its sources. Three primary
documents---the Quantification Guide, the operationalization addendum, and
PCBA Appendix~I---together with the PCBA engineering book's gain definitions,
describe what is nominally one framework. Read side by side, they do not
agree. They disagree about numeric scales, about how a full-scale reference is
obtained, about what the symbols $V$, $I$, and $R$ denote, about what the gain
$\beta$ is, about whether a mnemonic is a classification, about whether
network arithmetic is established or conjectural, about what a residual can
identify, and about where quantization band boundaries fall. A conventional
specification effort would have resolved these disagreements editorially:
pick the reading that seems most sensible, silently rewrite the others to
match, and present the result as the framework's single inherited meaning.
The CAS-Ohm audit refuses that move. Every disagreement found during source
inspection is preserved as a first-class, machine-readable evidence record in
\texttt{CAS\_OHM\_SOURCE\_CONFLICT\_EVIDENCE\_v0.1.json}
(schema \texttt{project\_coherence.cas\_ohm\_source\_conflict\_evidence.v0.1},
status \textsc{draft / open / not frozen}, as of 2026-08-04), and every
downstream artifact---the typed ontology of
Chapter~\ref{ch03:ontology} where present, the calibration profile of
Chapter~\ref{ch05:profile}, the response model of
Chapter~\ref{ch06:response}---is designed so that the conflict cannot be
papered over by construction.

This chapter reproduces that audit at evidence-record granularity. For each
of the eight registered conflicts we give the evidence records
verbatim-faithfully (source identifier, locator, and recorded fact), name the
conflict kind under which the record classifies the disagreement, analyze at
length why the disagreement matters to any downstream quantitative use, and
state the binding disposition and candidate rule that the CAS-Ohm 1.0 type
system enforces in response. We close with the record's own summary fields:
the conflict count, the gate effect, and the claim boundary---and with an
argument that the preservation itself, independent of any future empirical
outcome, is an engineering result.

A standing caution governs everything below. The audit record preserves
\emph{conflicts about texts}. It does not adjudicate which text is right,
because on the evidence available none of them is established: no source has
an admitted empirical validation, and the candidate as a whole is
\textsc{ring 4 / not licensed / held}. The dispositions quoted in this
chapter are candidate rules---typed constraints a successor implementation
enforces so that the ambiguity is representable and inert rather than
invisible and active. They are not verdicts on the sources.

\section{The four sources and their locator systems}\label{ch04:sources}

The evidence record opens with a source manifest. Four sources are
registered, each with a content digest, a declared locator system, and a
declared inspection method. We reproduce the manifest faithfully before
discussing why its structure matters.

\begin{description}
\item[\textsc{src-cas-ohm-guide-01}] (the Quantification Guide). SHA-256
\texttt{89AE13C0A5D6208E328E9EDC55F816609BB172076CC13093DCCCCF418145CE64}.
Locator system: \emph{PDF page number}. Inspection: relevant cited pages were
rendered and visually inspected.
\item[\textsc{src-cas-ohm-addendum-01}] (the operationalization addendum).
SHA-256
\texttt{0B62BBC66828A273E5458CC4D98EEE9B341554A474CCACB6387B209B5F0BE5E5}.
Locator system: \emph{\texttt{python-docx} paragraph and table-row index}.
Inspection: the OOXML content was inspected structurally. LibreOffice failed
with a PNG writer error and Word export did not complete, so \emph{no
page-number claim is made for this source}.
\item[\textsc{src-cas-ohm-pcba-app-i}] (PCBA Appendix~I). SHA-256
\texttt{9FFFAC3E16B43C6D809ECE913FD8C09E40C65F60F3B4BE693F45D35EC4E6AC45}.
Locator system: \emph{UTF-8 Markdown line number}. Inspection: the cited line
ranges were inspected directly from the source bytes.
\item[\textsc{src-pcba-engineer-20260726}] (the PCBA engineering book, cited
for its gain definitions). SHA-256
\texttt{294320B39FA521EAC086E99E2D9DB913E2B6CFF0754FF82DB1EF966424566274}.
Locator system: \emph{PDF page number and printed book page}. Inspection: the
cited gain-definition pages were rendered and visually inspected.
\end{description}

\subsection{Why disclosed locator systems matter}\label{ch04:locators}

A citation is a claim of the form ``at this position in that artifact, this
fact appears.'' The claim has three parts---artifact identity, position, and
fact---and each part can fail independently. The source manifest hardens all
three. Artifact identity is fixed by a SHA-256 content digest, so ``the
Guide'' is not a title that could silently drift across revisions but a
specific byte sequence. The fact is recorded verbatim in the evidence record,
as the following sections reproduce. The position, however, is the part most
often left informal in technical writing, and it is the part the manifest is
most careful about, because \emph{position is only meaningful relative to a
rendering procedure}, and different artifact formats admit different
renderings.

Consider each system in turn.

\emph{PDF pages, rendered and visually inspected.} A PDF fixes its own
pagination, so ``PDF page 6'' is well defined given the artifact bytes---but
only if someone actually rendered page 6 and looked at it. The inspection
note records that this was done. The distinction matters: a page number
harvested from a table of contents, a search index, or a prior citation is a
transcription, not an observation. The manifest asserts observation.

\emph{\texttt{python-docx} paragraph and table-row indices, with no page
number asserted.} A DOCX file has no intrinsic pagination at all. Its page
breaks are computed at render time by a layout engine, and different engines
(or the same engine with different fonts, margins, or versions) paginate
differently. The honest locator for a DOCX artifact inspected without a
successful rendering is therefore a \emph{structural} coordinate: the index
of a paragraph or a table row in the OOXML element sequence, which is a
property of the bytes themselves. The manifest goes further and records
\emph{why} no page number is asserted: both available page-rendering routes
failed---LibreOffice with a PNG writer error, and Word export did not
complete. This is a small but exemplary act of negative disclosure. The
auditors could have guessed a page number, or cited one from memory, and no
reader would likely have caught it. Instead the record states the failure and
confines its claims to coordinates it can defend. Every addendum citation in
this chapter therefore reads ``paragraph $n$'' or ``table row $k$,'' never
``page $p$.''

\emph{UTF-8 line numbers from exact bytes.} Appendix~I is a Markdown text
file. Line numbers in a text file are exact byte-level coordinates: given the
artifact digest, ``lines 20--31'' names a unique substring, checkable by
anyone holding the bytes, with no rendering step in between. This is the
strongest locator system of the four, and the record uses it wherever the
source format permits.

\emph{PDF and printed book pages together.} The PCBA engineering book exists
both as a PDF and as a printed volume whose pagination differs (front matter,
trim, and imposition shift the count). Citing both---``PDF page 56 (book page
40)''---lets a reader verify against either physical form without a
conversion table.

\begin{remark}\label{ch04:rem:locators}
The reason to belabor locator systems is that the entire audit inherits its
credibility from them. A conflict record asserts that sources say incompatible things --- or, for
the under-specification kinds, that a source's text fails to determine what
its own downstream use requires. That assertion is only as strong as the reader's ability
to open each source at the cited position and confirm the recorded fact. If
locators were informal---``somewhere in the Guide''---the conflicts would be
paraphrases, contestable and unfalsifiable. Because every locator is exact
relative to a disclosed rendering procedure (or, for the addendum, exact
relative to the bytes with the rendering honestly declared failed), each
conflict record is an independently checkable claim. The audit practices on
itself the evidential discipline the CAS-Ohm ontology imposes on
measurements: no bare assertions, provenance always attached, and missing
capabilities (here, DOCX pagination) recorded as missing rather than
interpolated.
\end{remark}

\subsection{Anatomy of a conflict record}\label{ch04:anatomy}

Each of the eight entries in the evidence record has the same shape, which
this chapter mirrors section by section.

\begin{definition}[Conflict record]\label{ch04:def:record}
A conflict record consists of: (i) a stable identifier (for example
\texttt{OHM-CONFLICT-SCALE-AND-SATURATION}); (ii) a \emph{kind}, drawn from a
small vocabulary describing the failure mode
(\texttt{DIRECT\_NUMERIC\_CONFLICT},
\texttt{INCOMPATIBLE\_CALIBRATION\_METHODS},
\texttt{UNIT\_UNDERSPECIFICATION},
\texttt{DISTINCT\_MATHEMATICAL\_PROFILES}, \texttt{MNEMONIC\_OVERLOAD},
\texttt{VALIDATION\_STATUS\_CONFLICT},
\texttt{CAUSAL\_OVERINTERPRETATION\_RISK},
\texttt{PRECISION\_UNDERSPECIFICATION}); (iii) an evidence list, each element
binding a source identifier from the manifest, a locator in that source's
declared locator system, and the recorded fact; (iv) a \emph{disposition},
naming the strategy the candidate adopts
(for example \texttt{PRESERVE\_AS\_PROFILE\_CONFLICT}); and (v) a
\emph{candidate rule}, the typed constraint a successor implementation
enforces.
\end{definition}

The kinds are worth a moment of attention because they are a taxonomy of how
multi-document specifications fail. A \texttt{DIRECT\_NUMERIC\_CONFLICT} is
the crudest failure: two sources assert incompatible numbers for the same
named thing. \texttt{UNIT\_UNDERSPECIFICATION} and
\texttt{PRECISION\_UNDERSPECIFICATION} are subtler: no single assertion is
wrong, but jointly the sources fail to determine a quantity's dimension or a
boundary's behavior. \texttt{INCOMPATIBLE\_CALIBRATION\_METHODS} and
\texttt{DISTINCT\_MATHEMATICAL\_PROFILES} mark places where one name covers
several genuinely different procedures or definitions.
\texttt{MNEMONIC\_OVERLOAD} marks a teaching device promoted into a
classification. \texttt{VALIDATION\_STATUS\_CONFLICT} marks disagreement not
about content but about epistemic standing---is a rule established or
conjectural? And \texttt{CAUSAL\_OVERINTERPRETATION\_RISK} marks language
that invites causal conclusions its own source elsewhere disclaims. Most
large specification corpora contain instances of most of these kinds; what is
unusual here is only that they are indexed.

Table~\ref{ch04:tab:overview} summarizes the eight records for orientation;
the sections that follow reproduce and analyze each in full.

\begin{table}[htbp]
\centering\small
\begin{tabular}{@{}lll@{}}
\toprule
Conflict (short name) & Kind & Disposition \\
\midrule
Scale and saturation & \texttt{DIRECT\_NUMERIC\_CONFLICT} &
\texttt{PRESERVE\_AS\_PROFILE\_CONFLICT} \\
Full-scale profile & \texttt{INCOMPATIBLE\_CALIBRATION\_METHODS} &
\texttt{SEPARATE\_NAMED\_PROFILES} \\
Dimensional bridge & \texttt{UNIT\_UNDERSPECIFICATION} &
\texttt{REQUIRE\_CALIBRATED\_UNIT\_BRIDGE} \\
Gain definition & \texttt{DISTINCT\_MATHEMATICAL\_PROFILES} &
\texttt{NO\_UNPROVED\_IDENTIFICATION} \\
ASOM axes & \texttt{MNEMONIC\_OVERLOAD} &
\texttt{SPLIT\_RESPONSE\_AND\_SEMANTIC\_AXES} \\
Network rules & \texttt{VALIDATION\_STATUS\_CONFLICT} &
\texttt{CAPABILITY\_GATE\_EACH\_OPERATION} \\
Residual identification & \texttt{CAUSAL\_OVERINTERPRETATION\_RISK} &
\texttt{NON\_IDENTIFYING\_DIAGNOSTIC\_ONLY} \\
Q4 boundaries & \texttt{PRECISION\_UNDERSPECIFICATION} &
\texttt{DEFINE\_EXACT\_HALF\_OPEN\_BINS} \\
\bottomrule
\end{tabular}
\caption{The eight preserved conflicts of
\texttt{CAS\_OHM\_SOURCE\_CONFLICT\_EVIDENCE\_v0.1.json}.}
\label{ch04:tab:overview}
\end{table}

\subsection{Verbatim fidelity}\label{ch04:fidelity}

One further discipline governs the presentation in this chapter. The
``fact'' field of each evidence record is a fixed string in the JSON
artifact, and this chapter reproduces those strings faithfully rather than
paraphrasing them. The distinction is not pedantry. A paraphrase of a
conflict is an interpretation, and interpretations of conflicting texts tend,
by small increments, toward whichever resolution the paraphraser privately
favors---which is exactly the harmonization pressure the audit exists to
resist. If the record says Appendix~I identifies $\mathrm{sat}=8.5$ ``as a
mixed-scale error whose component value is 0.85,'' then a restatement such
as ``Appendix~I corrects the Guide's threshold'' would already have taken a
side (that Appendix~I is the corrector rather than a second conflicting
voice), and a restatement such as ``the sources disagree slightly about
thresholds'' would have shrunk a category error to a rounding quibble.
Analysis in this chapter is therefore kept structurally separate from
evidence: the \emph{conflict} environments carry the recorded facts with
their locators, and the analysis subsections argue about consequences
without ever restating the facts in altered form. Readers auditing this
monograph against the JSON record should find every fact string, every
locator, every kind, every disposition, and every candidate rule matching
the artifact---and should treat any divergence they find as an error in this
text, not a permissible gloss.

\section{Conflict 1: Scale and Saturation}\label{ch04:scale}

\begin{conflict}[\texttt{OHM-CONFLICT-SCALE-AND-SATURATION}; kind
\texttt{DIRECT\_NUMERIC\_CONFLICT}]\label{ch04:conf:scale}
The evidence records are:
\begin{itemize}
\item \textsc{src-cas-ohm-guide-01}, PDF page 6: ``Drive is 0.0--10.0,
constraint $R$ is 0.1--10.0, and raw attention is computed as $V/R$.''
\item \textsc{src-cas-ohm-guide-01}, PDF page 12: ``The worked CJT uses
$\mathrm{sat}=8.5$ as a component threshold.''
\item \textsc{src-cas-ohm-pcba-app-i}, lines 20--31 and 51--59: ``Constraint
$R$ is 0--100, saturation/cutoff/threshold are 0--1, and $\mathrm{sat}=8.5$
is identified as a mixed-scale error whose component value is 0.85.''
\end{itemize}
Disposition: \texttt{PRESERVE\_AS\_PROFILE\_CONFLICT}.
\end{conflict}

\subsection{Analysis}

This is the audit's cleanest conflict and its most instructive, because the
disagreement is numeric and total. The Guide's page 6 places the constraint
quantity $R$ on a $0.1$--$10.0$ scale. Appendix~I's lines 20--31 place the
same-named quantity on a $0$--$100$ scale. These ranges are not related by a
stated conversion; they simply coexist under one symbol. Worse, the Guide's
worked example on page 12 uses $\mathrm{sat}=8.5$ as a component threshold,
while Appendix~I---in the very lines that fix saturation, cutoff, and
threshold quantities to the interval $0$--$1$---identifies that same $8.5$ as
a \emph{mixed-scale error} whose intended component value is $0.85$. One
source's worked pedagogy is another source's diagnosed bug.

Why does this matter beyond tidiness? Because every quantitative operation
downstream of these definitions is scale-sensitive. Consider the innocent
formula on Guide page 6: raw attention is $V/R$. If $R$ is drawn from
$0.1$--$10.0$, the quotient for a given $V$ differs by up to a factor of ten
from the quotient computed with $R$ drawn from $0$--$100$---and at the low
end of the Appendix~I scale, $R=0$ is admissible, making the quotient
undefined where the Guide's floor of $0.1$ guaranteed it finite. A threshold
comparison such as ``is saturation above $\mathrm{sat}$?'' flips its truth
value depending on whether the comparand lives in $[0,1]$ or $[0,10]$: a
fraction of $0.9$ is above a threshold of $0.85$ and far below a threshold of
$8.5$. A system that ingests both sources without typing the scales will,
with certainty rather than mere risk, compare quantities across scales and
produce numbers that are wrong by an order of magnitude while looking
perfectly well formed. Mixed-scale failure is silent: nothing about the
floating-point value $8.5$ announces which scale it came from.

The $\mathrm{sat}=8.5$ case is a documented instance of the failure actually
occurring \emph{inside the source corpus itself}. This is
not a hypothetical about what a careless integrator might do; the Guide's own
worked example, per Appendix~I's reading, already committed the mixed-scale
error that the framework's downstream users would be warned against. When a
specification's teaching material contains the bug its appendix diagnoses,
the specification cannot be said to carry one coherent scale system, and any
implementation claiming to inherit ``the'' scale system is choosing among
contradictory parents without saying so.

What the third evidence record does and does not assert must be stated
precisely. Appendix~I \emph{identifies} $\mathrm{sat}=8.5$ as a mixed-scale
error and states that its component value is $0.85$. That is a reading by one
source of another source's text---a diagnosis, with a proposed intended
value---not a repair the audit endorses. The audit's posture is deliberately
one level up: it records that the Guide asserts $8.5$, that Appendix~I calls
$8.5$ an error with intended value $0.85$, and that these positions
conflict. A harmonizing editor would have adopted Appendix~I's diagnosis,
rewritten the Guide's example, and moved on; the audit instead preserves
both, because adopting the diagnosis would presuppose exactly the
scale-priority question (which source's scale system governs?) that the
conflict exists to keep open. Unit-confusion failures are a storied genre in
engineering practice generally, and their common anatomy is visible here in
miniature: two communities of practice, each internally consistent, each
producing artifacts that interoperate numerically while diverging
semantically, with no type system standing between them. The lesson the
broader engineering literature draws---encode the unit in the type, not in
the documentation---is precisely the candidate rule's content.

\subsection{Disposition and candidate rule}

\begin{ruleenv}[Candidate rule for Conflict 1]\label{ch04:rule:scale}
``Q10 score, Q100 constraint score, normalized fraction, alpha dose,
activation projection, and effective resistance are distinct quantity types;
no silent conversion is allowed.''
\end{ruleenv}

The disposition \texttt{PRESERVE\_AS\_PROFILE\_CONFLICT} means the candidate
does not decide whether $R$ ``really'' lives on $0.1$--$10.0$ or $0$--$100$.
Instead the type system makes the question unaskable in ambiguous form: a
value carries \texttt{Q10} or \texttt{Q100} or \texttt{FRACTION} as its unit
kind, the quantity kind constrains which unit kinds are admissible, range
checks are enforced per unit ($[0,10]$ for \texttt{Q10}, $[0,100]$ for
\texttt{Q100}, $[0,1]$ for \texttt{FRACTION}), and no arithmetic path
converts between them without an explicit, provenance-carrying operation.
Under this rule the historical $\mathrm{sat}=8.5$ is not silently repaired to
$0.85$; it is \emph{unrepresentable} as a saturation fraction, because a
fraction-typed field rejects $8.5$ at validation. The conflict is preserved
in the record and neutralized in the types.

\section{Conflict 2: The Full-Scale Profile}\label{ch04:ifs}

\begin{conflict}[\texttt{OHM-CONFLICT-FULL-SCALE-PROFILE}; kind
\texttt{INCOMPATIBLE\_CALIBRATION\_METHODS}]\label{ch04:conf:ifs}
The evidence records are:
\begin{itemize}
\item \textsc{src-cas-ohm-guide-01}, PDF page 6: ``$I_{\mathrm{FS}}$ is a
selected board-specific full-scale attention value.''
\item \textsc{src-cas-ohm-addendum-01}, paragraph 30 and operationalization
table row 13: ``$I_{\mathrm{FS}}$ is the per-test-point 95th percentile of
absolute directional projection over a fixed hashed corpus.''
\item \textsc{src-cas-ohm-pcba-app-i}, lines 40--49 and 263--267:
``$I_{\mathrm{FS}}$ is measured once per path at drive rail and constraint
floor and does not transfer.''
\end{itemize}
Disposition: \texttt{SEPARATE\_NAMED\_PROFILES}.
\end{conflict}

\subsection{Analysis}

$I_{\mathrm{FS}}$ is the full-scale reference against which raw response is
normalized; every Q10 display value and every Q4 band assignment
(Conflict~\ref{ch04:conf:q4}) is a ratio with $I_{\mathrm{FS}}$ in the
denominator. The three sources give three procedures for obtaining it, and
the procedures differ in every property that matters for a denominator.

The Guide's $I_{\mathrm{FS}}$ is \emph{declared}: a board-specific full-scale
value that someone selected. It is a design decision, not a measurement; it
has no sampling distribution, no uncertainty, and no dependence on any
corpus. The addendum's $I_{\mathrm{FS}}$ is \emph{estimated}: the 95th
percentile of the absolute directional projection, computed per test point,
over a fixed hashed corpus. It is a statistic; it inherits sampling
variability from the corpus, changes if the corpus changes, and is pinned
down only because the corpus is content-addressed (``fixed hashed''). The
Appendix~I $I_{\mathrm{FS}}$ is \emph{probed at an extreme operating point}:
measured once per path with drive at its rail and constraint at its floor,
and---the source says so explicitly---it does not transfer.

These are not three descriptions of one number. In general, a declared design
ceiling, an empirical 95th percentile over ordinary traffic, and a maximal
excursion under rail/floor forcing will be three different magnitudes, and
nothing in any source relates them. The consequences propagate through every
normalized quantity. The same raw response yields different Q10 values under
the three denominators, hence potentially different Q4 bands, hence different
region classifications where thresholds are compared, hence different
downstream diagnostics. Two reports both stating ``$I_Q = 7.8$'' agree on
nothing unless they also agree on which full-scale procedure produced the
denominator---and under the corpus-percentile method, on which corpus, at
which percentile, and at which test point; and under the rail/floor method,
on which path, since that source warns the value does not transfer even
between paths.

The transferability warning is the sharpest edge. A declared full-scale might
plausibly be reused across a family of boards; a rail/floor probe is
explicitly local. Averaging the three, or substituting one when another is
unavailable, would produce a denominator with no defined meaning at all.

The three methods also differ in their failure modes, which is the deeper
reason they cannot be interchanged even ``approximately.'' A declared design
full-scale fails by becoming stale: the board changes, the declaration does
not, and nothing in the number signals the divergence---only a drift-check
discipline external to the value can catch it. A corpus percentile fails by
corpus shift and by tail sensitivity: the 95th percentile is an order
statistic of the projection distribution's upper tail, exactly the part of a
distribution most affected by a handful of outlying documents, so an
apparently minor corpus revision can move the denominator while every
individual document looks innocuous. A rail/floor probe fails by
extrapolation: it characterizes the system at an extreme operating point
that ordinary operation never visits, so its relevance to mid-range
normalization is an assumption, not an observation. Three methods, three
distinct threat models, three different maintenance obligations. A framework
that let one number stand for all three would not merely blur magnitudes; it
would blur which failure its operators needed to be watching for. Keeping
the methods separately named keeps their maintenance disciplines separately
owned.

\subsection{Disposition and candidate rule}

\begin{ruleenv}[Candidate rule for Conflict 2]\label{ch04:rule:ifs}
``Board-selected, corpus-percentile, and rail/floor full-scale procedures
remain different calibration methods with different profile digests.''
\end{ruleenv}

The disposition \texttt{SEPARATE\_NAMED\_PROFILES} is implemented as the
\texttt{IFSMethod} enumeration in the candidate code
(\texttt{RAIL\_AND\_FLOOR\_PER\_PATH},
\texttt{CORPUS\_PERCENTILE\_PER\_TEST\_POINT},
\texttt{DECLARED\_DESIGN\_FULL\_SCALE}), examined field by field in
Chapter~\ref{ch05:profile}. Because the method is a field of the calibration
profile and the profile's canonical digest covers it, two normalizations
performed under different full-scale methods are performed under different
profile digests, and the profile-relativity principle (nothing transfers
across profiles without a stated comparison rule) blocks their naive
comparison mechanically. No method is anointed correct; all three are
representable; none is substitutable for another.

\section{Conflict 3: The Dimensional Bridge}\label{ch04:bridge}

\begin{conflict}[\texttt{OHM-CONFLICT-DIMENSIONAL-BRIDGE}; kind
\texttt{UNIT\_UNDERSPECIFICATION}]\label{ch04:conf:bridge}
The evidence records are:
\begin{itemize}
\item \textsc{src-cas-ohm-guide-01}, PDF pages 6, 10, and 18: ``$V$, $R$, and
$I$ are design-scale quantities related by
$I_{\mathrm{pred}}=V/R_{\mathrm{total}}$.''
\item \textsc{src-cas-ohm-addendum-01}, operationalization table rows 1, 2,
5, and 12: ``$V$ is steering alpha, $I$ is directional activation projection,
contextual $R$ is inferred from a paired delta, and $I_{\mathrm{pred}}$ still
uses $V/R_{\mathrm{total}}$.''
\item \textsc{src-cas-ohm-pcba-app-i}, lines 90--108 and 326--339: ``The
relation is presented as testable design/calibration arithmetic, not as
physical voltage or a measured live-model law.''
\end{itemize}
Disposition: \texttt{REQUIRE\_CALIBRATED\_UNIT\_BRIDGE}.
\end{conflict}

\subsection{Analysis}

Between the Guide and the addendum, the symbols keep their algebra and change
their referents. On the Guide's pages, $V$, $R$, and $I$ are design-scale
scores---abstract magnitudes on the score scales of
Conflict~\ref{ch04:conf:scale}---and $I_{\mathrm{pred}}=V/R_{\mathrm{total}}$
is arithmetic among scores. In the addendum's operationalization rows, $V$
has become a steering alpha (an intervention magnitude applied to a model),
$I$ has become a directional activation projection (an observed quantity in a
model's representation space), and contextual $R$ is inferred from a paired
delta. Yet the prediction equation is carried over unchanged:
$I_{\mathrm{pred}}$ still uses $V/R_{\mathrm{total}}$.

Dimensional analysis makes the gap precise. For
$\hat I = V/R_{\mathrm{total}}$ to be dimensionally coherent with $V$ in
alpha units and $I$ in activation units, $R_{\mathrm{total}}$ must carry
units of alpha per activation \cite{buckingham}. But the $R$ components the
sources supply are score-derived---dimensionless design-scale numbers from
Conflict~\ref{ch04:conf:scale}'s (mutually conflicting) ranges. Dividing an
alpha by a dimensionless score yields alpha, not activation: the formula, as
inherited, predicts a quantity in the wrong dimension. Nothing in any source
supplies the missing conversion---the calibrated bridge that would turn a
score-derived resistance into a quantity with alpha-per-activation units.
This is why the record's kind is \texttt{UNIT\_UNDERSPECIFICATION} rather
than a direct conflict: no single source asserts a falsehood; jointly they
fail to define the one constant that would make the central equation mean
anything.

Appendix~I's contribution is a guard on interpretation: the relation is
testable design/calibration arithmetic, not physical voltage and not a
measured live-model law. This matters in both directions. It blocks the
physical reading (nothing here is electrically a volt, an ampere, or an ohm;
the candidate's unit ontology makes those forbidden aliases), and it blocks
the empirical reading (writing the equation does not make it a law of any
live model; whether any $(a, y)$ pair from a real system satisfies a linear
relation is exactly what the registered hypotheses of
Chapter~\ref{hyp:chapter} where present leave to controlled experiments).

The double negation in Appendix~I's framing---not physical, not yet
empirical---leaves the equation in a precisely intermediate status that the
candidate's vocabulary makes explicit: it is \emph{design arithmetic}, a
relation among declared and calibrated quantities that is testable in
principle once its units are fixed, and nothing more until tested. The
physical exclusion is enforced in the unit ontology as a forbidden-alias
rule: \textsc{volt}, \textsc{ampere}, and \textsc{ohm} are not merely absent
from the unit enumeration but affirmatively barred as labels for model
quantities, so that the analogical vocabulary of the sources ($V$, $I$, $R$)
can never quietly reify into a claim of electrical measurement. The
empirical exclusion is enforced by the response model's admissibility
conditions (Chapter~\ref{ch06:response}): the equation may be evaluated only
with a calibrated, uncertainty-carrying, held-out-validated effective
resistance, which as of this writing does not exist for any live system.

The practical stakes are high because $\hat I = V/R_{\mathrm{total}}$ is the
framework's load-bearing equation: it is the prediction whose residuals drive
the diagnostics of Conflict~\ref{ch04:conf:residual} and whose network
compositions drive Conflict~\ref{ch04:conf:networks}. An equation with an
undefined dimension in its denominator contaminates everything built on it.
Any residual computed from such a prediction is a difference between an
observed activation and a quantity of indeterminate dimension; any series sum
of such resistances aggregates numbers whose common unit was never
established.

\subsection{Disposition and candidate rule}

\begin{ruleenv}[Candidate rule for Conflict 3]\label{ch04:rule:bridge}
``Effective resistance may be used only with an explicit alpha-unit per
activation-unit calibration; constraint scores never masquerade as that
quantity.''
\end{ruleenv}

Under \texttt{REQUIRE\_CALIBRATED\_UNIT\_BRIDGE}, the candidate splits the
overloaded $R$ into types that cannot be confused: \texttt{CONSTRAINT\_SCORE}
(unit \texttt{Q10} or \texttt{Q100}) for the score-scale quantity, and
\texttt{EFFECTIVE\_RESISTANCE} (unit \texttt{ALPHA\_PER\_ACTIVATION},
strictly positive) for the quantity admissible in the prediction
$\hat y = a/r_{\mathrm{eff}}$. There is no constructor from one to the other.
An effective resistance comes into existence only through a declared fit on a
calibration split under a full calibration profile, carrying uncertainty and
a role tag (Chapter~\ref{ch06:response}); a constraint score can never be
passed where an effective resistance is required, because the unit kinds
differ and unit checking is mandatory. The Guide's equation shape is thereby
preserved as a \emph{hypothesis about a calibrated quantity} while the
dimensional hole that made its inherited form unmeaning is closed by
type-level refusal rather than by silent reinterpretation.

\section{Conflict 4: The Gain Definition}\label{ch04:gain}

\begin{conflict}[\texttt{OHM-CONFLICT-GAIN-DEFINITION}; kind
\texttt{DISTINCT\_MATHEMATICAL\_PROFILES}]\label{ch04:conf:gain}
The evidence records are:
\begin{itemize}
\item \textsc{src-pcba-engineer-20260726}, PDF page 56 (book page 40) and PDF
page 153 (book page 137): ``CJT beta is defined by Fisher--Rao small- or
large-signal displacement ratios on the open simplex.''
\item \textsc{src-cas-ohm-addendum-01}, operationalization table row 6:
``Beta is operationalized as a local Jacobian/projection derivative between
cause and effect test points.''
\item \textsc{src-cas-ohm-guide-01}, PDF page 13: ``Stage beta values
multiply in cascades while attention flow is local at each stage.''
\end{itemize}
Disposition: \texttt{NO\_UNPROVED\_IDENTIFICATION}.
\end{conflict}

\subsection{Analysis}

The symbol $\beta$ names a gain---an amplification factor between a cause
stage and an effect stage. The PCBA engineering book defines it
geometrically: a ratio of Fisher--Rao displacements on the open probability
simplex, in small-signal or large-signal form. The addendum operationalizes
it analytically: a local Jacobian/projection derivative between cause and
effect test points---a directional sensitivity of one projection with respect
to another in a model's representation space. The Guide, meanwhile, uses
$\beta$ compositionally: stage values multiply along cascades, while
attention flow remains local at each stage.

The Fisher--Rao definition and the Jacobian/projection definition are
different mathematical objects in almost every respect that could matter.
They live on different spaces: the open simplex with its information
geometry, versus activation space with a Euclidean projection structure. They
have different invariance properties: Fisher--Rao displacement is invariant
under the simplex's natural reparameterizations, while a projection
derivative depends on the chosen direction vectors at both test points. They
respond differently to signal size: the book itself distinguishes small- and
large-signal forms, while a local derivative is by construction a
small-signal object. There is no theorem in any of the sources---and the
record is explicit that none is currently supplied---identifying the two, nor
even a coordinate map under which such an identification could be stated.

The Guide's cascade rule then quietly raises the stakes. Multiplying stage
gains along a cascade is exactly the operation that is \emph{most} sensitive
to which definition is in force, because composition of local derivatives
follows a chain rule only when adjacent stages' coordinates match, and
composition of Fisher--Rao displacement ratios has no chain rule at all
without further structure. Under the multiplication rule, an error of
identification does not stay local: it compounds geometrically with cascade
depth. Three stages of a modest 30\% definitional discrepancy are a factor of
$2.2$; the deeper the cascade, the worse. (Those numbers are illustrative
arithmetic about error propagation, not measurements.)

The general lesson is about symbol overloading across document genres. A
definitional text (the book), an operationalization (the addendum), and a
usage idiom (the Guide's cascades) each pulled $\beta$ toward their own
needs. Each is internally reasonable. Jointly they leave ``the gain''
denoting two unbridged mathematical profiles plus a composition rule that
presupposes the bridge.

\subsection{Disposition and candidate rule}

\begin{ruleenv}[Candidate rule for Conflict 4]\label{ch04:rule:gain}
``Fisher--Rao gain and local Jacobian/projection gain are separately named
profiles until a theorem or experiment supplies a coordinate bridge.''
\end{ruleenv}

\texttt{NO\_UNPROVED\_IDENTIFICATION} is the mathematical analogue of
Conflict~\ref{ch04:conf:ifs}'s disposition: one name, several referents,
therefore several names. The candidate types gain magnitudes as dimensionless
quantities bound to a declared method identity, so a Fisher--Rao displacement
gain and a Jacobian/projection gain carry different method identifiers and
cannot be pooled, averaged, or multiplied together by accident. The gain
composition operation (Chapter~\ref{ch09:networks} where present) is
capability-gated and requires adjacent coordinate agreement precisely so that
the Guide's cascade multiplication can be \emph{tested} as the registered
hypothesis \textsc{clm-cas-ohm-gain-composition} rather than assumed. If a
future theorem or a controlled experiment supplies the bridge, the
identification becomes a recorded, evidenced event; until then the two
profiles remain separate by name and by type.

\section{Conflict 5: ASOM Axis Overload}\label{ch04:asom}

\begin{conflict}[\texttt{OHM-CONFLICT-ASOM-AXES}; kind
\texttt{MNEMONIC\_OVERLOAD}]\label{ch04:conf:asom}
The evidence records are:
\begin{itemize}
\item \textsc{src-cas-ohm-guide-01}, PDF page 11: ``ASOM is explicitly
session-derived and overlays Starved/Aligned/Oversaturated/Misaligned onto
cutoff/active/saturation/breakdown.''
\item \textsc{src-cas-ohm-addendum-01}, paragraphs 14--15 and ASOM table rows
1--4: ``The addendum uses dose-response signatures plus effect-axis outcomes
and warns that flat response has multiple causes.''
\end{itemize}
Disposition: \texttt{SPLIT\_RESPONSE\_AND\_SEMANTIC\_AXES}.
\end{conflict}

\subsection{Analysis}

ASOM---Starved, Aligned, Oversaturated, Misaligned---is, per the Guide's own
page 11, session-derived: a mnemonic overlay mapping four evocative words
onto the four dose-response regions cutoff, active, saturation, and
breakdown. The overlay is pedagogically appealing because it compresses two
questions into one word. But the two questions are different in kind.
\emph{Where on the dose-response curve is the system operating?} is a
question about response shape: is response flat at low dose (cutoff), rising
(active), flat at high dose (saturation), anomalous (breakdown)?
\emph{Is the effect the intended one?} is a question about semantic outcome:
did the intervention move the targeted quantity in the targeted way, the
opposite way, or off-axis? The mnemonic fuses them: ``Aligned'' answers the
second question with the vocabulary slot of the first, and ``Misaligned''
imports a semantic verdict into what is structurally a region label.

The addendum resists the fusion. Its paragraphs 14--15 and ASOM table rows
1--4 keep dose-response signatures and effect-axis outcomes as separate
columns, and issue the warning that unravels the overlay: flat response has
multiple causes. A flat dose-response curve is compatible with cutoff (dose
below threshold), with saturation (dose beyond the linear region), with a
wrong direction vector (the probe measures an axis the intervention does not
move), with genuine insensitivity, and with instrument failure. A mnemonic
that reads ``flat'' as ``Starved'' has silently chosen one cause out of five.
Symmetrically, an active region says nothing about intent: a system can
respond briskly and in exactly the wrong direction, which the fused
vocabulary can only render as a paradox (``actively misaligned'') rather than
as the ordinary factored state (region \textsc{active}, outcome
\textsc{opposite}).

There is also an honest-ignorance problem the mnemonic cannot express. Four
words partition every observation into four verdicts; there is no ASOM cell
for ``the response region is known but the semantic outcome has not been
assessed,'' although that is the \emph{normal} state of a fresh measurement:
region classification needs only the dose-response sweep, while outcome
classification needs an independent determination of what the intervention
was supposed to do and whether it did it. A vocabulary without an
unassessed value forces premature verdicts by construction---every
observation must land somewhere, so it lands in whichever verdict its shape
most resembles. The two-axis scheme, by contrast, admits
\texttt{UNASSESSED} on each axis independently, so a system can honestly
report ``region \textsc{active}, outcome \textsc{unassessed}'' and
downstream consumers can distinguish ``verified as intended'' from ``not
yet examined''---a distinction on which, for a framework aimed at auditing,
essentially everything depends.

The failure kind, \texttt{MNEMONIC\_OVERLOAD}, names a pattern endemic to
maturing frameworks: a teaching aid, introduced honestly as such
(``session-derived''), hardens through repetition into an analytical
category, and downstream text begins drawing inferences from the mnemonic's
connotations rather than from the measurements the mnemonic was compressing.
The cost is not hypothetical. A classifier emitting ``Misaligned'' as a
single-axis verdict makes a semantic claim on the strength of a response
shape; on a system under audit, that is precisely the kind of unearned
conclusion the whole evidence discipline exists to prevent.

\subsection{Disposition and candidate rule}

\begin{ruleenv}[Candidate rule for Conflict 5]\label{ch04:rule:asom}
``Dose-response region and semantic outcome are classified independently;
active does not imply intended and non-ohmic does not imply misaligned.''
\end{ruleenv}

Under \texttt{SPLIT\_RESPONSE\_AND\_SEMANTIC\_AXES}, the candidate carries
two enumerations: dose-response region (\texttt{UNASSESSED},
\texttt{CUTOFF}, \texttt{ACTIVE}, \texttt{SATURATION}, \texttt{BREAKDOWN},
\texttt{NON\_OHMIC}, \texttt{UNSTABLE}---the source's breakdown region
retained) and semantic outcome (\texttt{UNASSESSED}, \texttt{INTENDED},
\texttt{OPPOSITE}, \texttt{OFF\_AXIS}). Every classification is a pair; no
region value mechanically determines an outcome value. ASOM survives---but
demoted to what the Guide said it was: a derived display label computed from
the pair after the fact, never an input, never an axis, and undefined
whenever the pair does not fit its four cells. The mnemonic keeps its
pedagogy and loses its authority.

\section{Conflict 6: Network Rules}\label{ch04:networks}

\begin{conflict}[\texttt{OHM-CONFLICT-NETWORK-RULES}; kind
\texttt{VALIDATION\_STATUS\_CONFLICT}]\label{ch04:conf:networks}
The evidence records are:
\begin{itemize}
\item \textsc{src-cas-ohm-guide-01}, PDF pages 13--14 and 18: ``Series
addition, reciprocal parallel aggregation, and beta multiplication are
presented as working network rules.''
\item \textsc{src-cas-ohm-pcba-app-i}, lines 206--254: ``The same arithmetic
is applied to paths, with a separate warning that co-present guardrails are
not parallel routes and flow remains local.''
\item \textsc{src-cas-ohm-addendum-01}, paragraphs 16--19 and falsification
table rows P2--P3: ``Gain composition and parallel conductance are
preregistered predictions whose rules must be dropped if the tests fail.''
\end{itemize}
Disposition: \texttt{CAPABILITY\_GATE\_EACH\_OPERATION}.
\end{conflict}

\subsection{Analysis}

Here the three sources agree on the arithmetic and disagree on its standing.
The rules in question are the circuit-theory triad: series resistances add;
parallel resistances aggregate by reciprocals; stage gains multiply. The
Guide presents them as \emph{working rules}---usable machinery. Appendix~I
applies the same arithmetic to paths but attaches a structural warning:
co-present guardrails are not parallel routes, and flow remains local. The
addendum reclassifies the rules entirely: gain composition and parallel
conductance are \emph{preregistered predictions}, indexed as falsification
rows P2 and P3, whose rules must be dropped if the tests fail.

A \texttt{VALIDATION\_STATUS\_CONFLICT} is in some ways more dangerous than a
numeric conflict, because the disputed content circulates freely. A reader of
the Guide alone will compute series equivalents and cascade gains as
routinely as an electrical engineer computes them for resistor networks---for
resistors, the rules are consequences of established physics. Here they are
borrowed shape. Whether ablating one of two ``parallel'' safety mechanisms
changes system response according to reciprocal aggregation is an empirical
question about a particular system, and the addendum has the epistemically
correct posture: state the prediction, register the test, and commit in
advance to dropping the rule on failure.

Appendix~I's warning identifies the specific modeling error most likely to
produce false parallel structure. Two guardrails that are co-present---both
active on the same trajectory---are not thereby alternative routes between
the same endpoints. Parallel aggregation presupposes genuine alternatives:
distinct branches with common endpoints such that flow divides among them.
Co-presence is a weaker, easier-to-observe relation, and conflating it with
parallelism imports an aggregation formula into a topology that does not
support it. The further warning that ``flow remains local'' blocks the
complementary error at the series end: chaining stage arithmetic into a
claim about a global conserved flow through the chain, when nothing
establishes a conservation law for attention flow across stages.

The record therefore preserves a three-way disagreement about \emph{how much
is known}: working rule (Guide), conditionally applicable arithmetic with
topological caveats (Appendix~I), falsifiable conjecture (addendum). An
implementation must choose a posture, and the choice is consequential: treat
conjecture as machinery and every downstream aggregate silently inherits an
unvalidated assumption; treat machinery as forbidden and legitimate
structural exploration is blocked.

The capability-gate disposition is best understood as a third posture that
dominates both extremes. Against the permissive posture, it ensures that no
aggregate can exist without a machine-checked record of the preconditions it
satisfied and---equally important---of the epistemic label it carries:
every network result in the candidate is stamped as an unauthenticated
structural research calculation, bound to its trust mode, its authority
manifest identity, and its declared covariance limitation, so the result
cannot later be quoted shorn of its status. Against the prohibitive posture,
it preserves the ability to \emph{compute} the structural quantities the
registered hypotheses will need: testing whether parallel conductance
predicts ablation response (hypothesis row P3) requires computing the
reciprocal aggregate to have a prediction to test, and a framework that
forbade the arithmetic outright could never run its own falsification
program. The gate thus operationalizes the addendum's preregistration
posture in code: the arithmetic runs, its outputs are labeled conjectural,
and the label can only be removed by the registered experiments---never by
repetition or lapse of memory.

\subsection{Disposition and candidate rule}

\begin{ruleenv}[Candidate rule for Conflict 6]\label{ch04:rule:networks}
``No series, parallel, or gain composition is enabled without coordinate,
unit, topology, active-region, and evidence compatibility under the exact
profile.''
\end{ruleenv}

\texttt{CAPABILITY\_GATE\_EACH\_OPERATION} threads the needle. The
operations exist in the candidate---but only as unauthenticated structural
research calculations, each behind a gate that checks precisely the
preconditions the sources' caveats identify: compatible coordinates (against
the false-parallelism error), a single alpha-per-activation unit (against
Conflict~\ref{ch04:conf:bridge}), genuine acyclic topology with distinct
endpoints (against co-presence masquerading as parallelism), active-region
components with adequate linearity (against composing outside the region
where a linear rule could even be a candidate), and complete evidence under
one exact profile. Every aggregate binds its child order, topology, and the
declared covariance limitation, and duplicate or cyclic paths fail closed.
The addendum's falsification posture is preserved intact: rows P2--P3
correspond to registered hypotheses, and the gates ensure that structural
results produced meanwhile are labeled as exactly that---structure, not
validated network law.

\section{Conflict 7: Residual Identification}\label{ch04:residual}

\begin{conflict}[\texttt{OHM-CONFLICT-RESIDUAL-IDENTIFICATION}; kind
\texttt{CAUSAL\_OVERINTERPRETATION\_RISK}]\label{ch04:conf:residual}
The evidence records are:
\begin{itemize}
\item \textsc{src-cas-ohm-guide-01}, PDF pages 10 and 16: ``Residual sign is
used to suggest hidden opposition, extra amplification, or hidden resistance,
while the page also says residual alone is insufficient.''
\item \textsc{src-cas-ohm-pcba-app-i}, lines 171--202 and 260--277:
``$R_{\mathrm{hidden}}$ is called a residual rather than a component, but
positive residual on a safety board is treated as a bypass and
saturated-stage residuals are meaningless.''
\item \textsc{src-cas-ohm-addendum-01}, paragraph 19 and falsification table
row P4: ``Hidden content is last in the diagnostic walk and residual validity
requires independently identified causes above a shuffled-control baseline.''
\end{itemize}
Disposition: \texttt{NON\_IDENTIFYING\_DIAGNOSTIC\_ONLY}.
\end{conflict}

\subsection{Analysis}

A residual is a subtraction: observed response minus predicted response. The
temptation the record documents is the promotion of that subtraction into an
identification---reading the sign of $\delta$ as the presence of a named
hidden mechanism. The Guide exhibits the temptation and its own antidote on
the same pages: residual sign ``suggests'' hidden opposition, extra
amplification, or hidden resistance, while the same page concedes that
residual alone is insufficient. Appendix~I shows the slide completed in two
steps: first terminological ($R_{\mathrm{hidden}}$ is ``called a residual
rather than a component''---the name of a model-mismatch symptom is dressed
as a circuit element), then inferential (a positive residual \emph{on a
safety board} is treated as a bypass). It simultaneously notes a condition
under which residuals carry no information at all: at a saturated stage,
where predicted response is flat, residuals are meaningless.

Why is the identification invalid? Because the residual is computed against a
fitted model, and every deficiency of that model lands in the residual
indistinguishably. A nonzero $\delta$ is compatible with: a genuine
unmodeled mechanism; misfit of $r_{\mathrm{eff}}$ (wrong region, drifted
calibration); a wrong or degraded direction vector; ordinary noise at the
recorded uncertainty; a stale profile; or an error in the observation chain.
Sign does not separate these. Inferring ``bypass'' from ``positive residual
on a safety board'' is an inference to the most alarming member of an
equivalence class the data cannot resolve---and its cost profile is
asymmetric in the worst way, since a bypass claim about a safety mechanism
triggers exactly the kind of downstream action that should require the
\emph{most} evidence, on the strength of a statistic that here carries the
least. The general point is Pearl's: causal identification requires an
explicit causal model and identifiable interventions, not signs of
differences from a curve fit \cite{pearl}.

The addendum again supplies the disciplined posture, in two parts. Hidden
content is \emph{last} in the diagnostic walk: before entertaining a hidden
mechanism one must exhaust the identifiable causes---region, calibration,
direction, noise, staleness. And residual validity is itself an empirical
claim with a registered test (falsification row P4): residual diagnostics
must localize independently identified causes above a shuffled-control
baseline. If shuffled controls localize as well as real residuals, the
diagnostic is noise and its causal vocabulary must be withdrawn.

The ``diagnostic walk'' ordering encodes a cost asymmetry that pure
statistics does not capture. The identifiable
causes ahead of hidden content in the walk---wrong region, stale
calibration, degraded direction, noise---are all \emph{checkable} at bounded
cost: each has an associated record in the profile or the observation chain
that can be examined directly. ``Hidden mechanism,'' by contrast, is
unbounded: it names the complement of everything modeled, so no finite check
can confirm it and almost any anomaly can be retrofitted to it. Placing it
last is therefore not a mere tidiness preference; it is the only ordering
under which the diagnostic process terminates against checkable facts before
reaching the unfalsifiable residue. Inverting the walk---reaching for the
hidden mechanism first, as the ``bypass'' reading does---converts every
calibration bug and every noisy run into apparent evidence of concealed
behavior, which for a safety-relevant audit is the most corrosive possible
failure: it manufactures alarming findings at exactly the rate of the
auditor's own instrumentation defects.

\subsection{Disposition and candidate rule}

\begin{ruleenv}[Candidate rule for Conflict 7]\label{ch04:rule:residual}
``Residuals retain propagated uncertainty and may rank follow-up probes; they
never prove a hidden policy, bypass, or cause. Zero or zero-compatible
observations are censored.''
\end{ruleenv}

\texttt{NON\_IDENTIFYING\_DIAGNOSTIC\_ONLY} assigns residuals their honest
job. In the candidate, a residual exists only for an admitted active-region
prediction, with instrumented held-out observations, matching test points and
profile, calibrated resistance, and propagated uncertainty; at zero or
zero-compatible observed response no ratio quantity is emitted at all (the
observation is censored with a recorded reason rather than divided into a
meaningless ratio), and the saturated-stage case Appendix~I flags is
excluded upstream by the active-region admission gate. What a
residual may do is \emph{rank follow-up probes}: allocate investigative
attention. What it may never do is appear as the proof of a hidden policy,
bypass, opposition, or cause. The registered hypothesis
\textsc{clm-cas-ohm-residual-localization} keeps even the ranking use on
probation: blinded implanted faults, matched controls, and withdrawal of
causal language on failure.

\section{Conflict 8: Q4 Boundaries}\label{ch04:q4}

\begin{conflict}[\texttt{OHM-CONFLICT-Q4-BOUNDARIES}; kind
\texttt{PRECISION\_UNDERSPECIFICATION}]\label{ch04:conf:q4}
The evidence record is:
\begin{itemize}
\item \textsc{src-cas-ohm-guide-01}, PDF page 7: ``Q4 is displayed with
two-decimal bands 0.00--2.49, 2.50--4.99, 5.00--7.49, and 7.50--10.00 without
a rule for higher-precision boundary values.''
\end{itemize}
Disposition: \texttt{DEFINE\_EXACT\_HALF\_OPEN\_BINS}.
\end{conflict}

\subsection{Analysis}

This is the audit's smallest conflict and the easiest to underestimate. The
Guide's Q4 quantizer displays four bands with two-decimal endpoints:
0.00--2.49, 2.50--4.99, 5.00--7.49, 7.50--10.00. As a display convention for
two-decimal values this is complete. As a specification of a function on
finer-grained inputs it is not, because the bands as written do not cover the
line: the value $2.495$ falls between 2.49 and 2.50, in no band. What happens
to it depends on an unstated rounding rule---and the two natural
resolutions disagree. Round-then-bin sends $2.495$ to $2.50$ (round-half-up)
and thence to the second band; bin-by-truncation keeps it under $2.50$ in the
first. The specification gap is only two decimal places wide, but a
classification function undefined on a set of inputs, with plausible readings
that disagree on those inputs, is not one function; it is a family of
functions wearing one name.

The failure has an aggravating property the other conflicts lack: it sits
exactly where decisions concentrate. Band boundaries are where a quantizer's
output changes, so any downstream logic keyed to bands---reporting tiers,
threshold alarms, comparisons across runs---changes behavior precisely at the
values the specification leaves undefined. Systems accumulate values near
their thresholds (thresholds are typically set where the interesting mass
is), so the undefined set is visited in practice, not just in principle.
Worse, an implementation that rounds before binning couples the
classification to display precision: increase the displayed decimals and
boundary values silently change bands, so the same measurement classifies
differently depending on a formatting choice. A quantization band should be a
fact about a value, not about its printout.

\subsection{Disposition and candidate rule}

\begin{ruleenv}[Candidate rule for Conflict 8]\label{ch04:rule:q4}
``The successor uses exact unrounded comparison against 2.5, 5, and 7.5;
normalization may not round a value across a band boundary.''
\end{ruleenv}

\texttt{DEFINE\_EXACT\_HALF\_OPEN\_BINS} completes the specification with
lower-closed half-open intervals. The mirror convention --- upper-closed
bins $(0,2.5]$, $(2.5,5]$, and so on --- would be equally total and equally
exact, but it would assign $2.50$ to the first band, contradicting the
Guide's own displayed membership, in which $2.50$ opens the band
``2.50--4.99''; the lower-closed convention is therefore the unique total
completion consistent with the source's display. The chosen intervals are
$[0,2.5)$, $[2.5,5)$,
$[5,7.5)$, $[7.5,\infty)$ evaluated on the exact pre-rounding value, with
boundary points assigned to the upper band by the half-open convention. The
candidate implements the comparison over exact rational arithmetic on the
retained pre-clipping ratio, so no representable input is ambiguous and
display rounding, which happens afterward and never re-enters
classification, cannot move a value across a boundary
(Chapter~\ref{ch08:quantization} where present states this as a
proposition). The Guide's two-decimal table remains exactly what it was---a
correct display of the bands at two decimals---while the function it
under-specified is now total.

\section{Interactions Among the Conflicts}\label{ch04:interactions}

The eight records are registered independently, but they are not independent
in effect. Reading them as a system reveals a dependency structure that
explains why the candidate's dispositions had to be typed constraints rather
than local patches: several conflicts compound, and a resolution of one that
ignored its neighbors would leak.

\emph{Conflict 1 feeds Conflict 3.} The dimensional bridge problem is
aggravated by the scale problem beneath it. Even if one resolved
Conflict~\ref{ch04:conf:bridge} by fiat---declaring that $R_{\mathrm{total}}$
shall carry alpha-per-activation units obtained from some conversion of
constraint scores---the conversion would have to know which score scale it
was converting \emph{from}, and Conflict~\ref{ch04:conf:scale} establishes
that the sources disagree about exactly that. A bridge from an ambiguous
abutment is ambiguous end to end: the same nominal constraint value of, say,
$5$ would bridge to resistances differing by an order of magnitude depending
on whether it was read on the $0.1$--$10.0$ or the $0$--$100$ scale. The
candidate rule for Conflict 1 (distinct quantity types, no silent
conversion) is therefore a load-bearing prerequisite for the candidate rule
of Conflict 3 (explicit calibrated bridge only): the bridge requirement is
meaningful because the scales it would bridge from are typed.

\emph{Conflict 2 feeds Conflict 8.} Q4 band assignment is a comparison of a
normalized ratio against fixed cut points, and the ratio's denominator is
$I_{\mathrm{FS}}$. Conflict~\ref{ch04:conf:ifs} means the denominator is
method-relative; Conflict~\ref{ch04:conf:q4} means the comparison was
boundary-underspecified. The two compound multiplicatively: a value near a
band boundary under one full-scale method may sit far from any boundary
under another, so the boundary rule's importance is itself method-dependent,
and no statement of the form ``this response is Q3'' is complete without
both the exact bin convention (Conflict 8's disposition) and the full-scale
method with its parameters (Conflict 2's disposition) attached. This is why
the candidate binds the \texttt{IFSMethod} into the profile digest and
computes bands from exact pre-rounding ratios: each disposition closes one
half of a two-part ambiguity in every quantized report.

\emph{Conflict 3 underlies Conflicts 6 and 7.} Series addition and
reciprocal parallel aggregation are operations \emph{on resistances}; the
residual is a difference between an observation and a prediction \emph{made
with a resistance}. Until the dimensional bridge exists, the operands of
Conflict~\ref{ch04:conf:networks}'s arithmetic and the predictions behind
Conflict~\ref{ch04:conf:residual}'s diagnostics are quantities of undefined
dimension. The capability gates on network operations therefore check unit
kind before topology---an aggregate of score-typed pseudo-resistances is
rejected at the unit check, before the question of whether its branches are
genuinely parallel even arises. Likewise the residual pipeline requires a
calibrated resistance with the correct unit as a precondition for emitting
any residual at all. Had the audit ``resolved'' Conflict 3 by tolerating
score-derived resistances, the gates of Conflicts 6 and 7 would be guarding
operations whose operands were already meaningless.

\emph{Conflict 4 arms Conflict 6.} The Guide's cascade rule multiplies stage
gains; Conflict~\ref{ch04:conf:gain} shows ``gain'' has two unbridged
definitions. The composition gate's requirement of adjacent coordinate
agreement and a single declared gain method is what prevents a cascade from
mixing Fisher--Rao displacement ratios with Jacobian derivatives---a mixture
that would be numerically computable (both are dimensionless) and
semantically void. This is the general pattern: dimensionless quantities are
the most dangerous operands in the system, because unit checking alone
cannot distinguish them, so method identity must travel with the value.

\emph{Conflict 5 disciplines Conflict 7's vocabulary.} The residual
temptation (``positive residual on a safety board is a bypass'') and the
mnemonic temptation (``flat response is Starved'') are the same move at
different altitudes: collapsing a measurement-shaped observation into a
semantic verdict without the evidence that would distinguish the verdict
from its alternatives. The split-axes disposition gives the system a place
to put honest ignorance---semantic outcome \texttt{UNASSESSED}---and the
residual disposition uses the same discipline: a residual may flag
\emph{that} something is unexplained, never \emph{what} explains it. Both
dispositions route the missing determination to the same destination:
independent evidence under registered protocols, not connotation.

\begin{finding}\label{ch04:find:web}
The eight conflicts form a dependency web in which the scale and bridge
conflicts (1 and 3) are upstream of nearly everything quantitative, the
method-identity conflicts (2 and 4) are upstream of all cross-context
comparison and composition, and the interpretation conflicts (5 and 7)
govern what may be said about the results. No single conflict could have
been resolved in isolation without either presupposing a resolution of its
upstream neighbors or silently exporting its ambiguity downstream. This is
the structural justification for the audit's uniform strategy: preserve
every conflict, and encode every disposition at the type level, where the
dependencies are enforced in composition order automatically.
\end{finding}

\section{The Record as a Whole}\label{ch04:close}

The evidence record closes with summary fields that are as much a part of the
audit as the conflicts themselves, and we reproduce them faithfully.

\emph{Conflict count: 8.} The count is a completeness assertion with a
precise scope: these eight records are what the inspection of the cited
sources, at the cited locators, under the declared inspection methods,
registered as conflicts or under-specifications. It is not an assertion that
the sources contain exactly eight problems; it is an assertion that the audit
found and preserved these eight, and that downstream artifacts claiming to
rest on the audit inherit exactly this list---no silent additions, no silent
deletions.

\emph{Gate effect: \texttt{KEEP\_OPEN}.} The audit feeds a gated development
process, and a source audit that discovered eight unresolved conflicts could
not honestly close anything. \texttt{KEEP\_OPEN} records that the audit's
effect on its gate is to hold it open: the conflicts are preserved, not
resolved; the registered hypotheses that would resolve the resolvable ones
(network rules, residual validity, gain composition) remain untested; and no
downstream claim may cite the audit as if it had settled what it explicitly
left unsettled.

\emph{The claim boundary.} The record's own boundary statement reads: ``This
record preserves source conflicts and under-specifications. It does not
validate CAS-Ohm empirically, identify physical quantities, license a
product claim, or authorize model intervention.'' The four exclusions are
worth reading as the four ways an audit is most easily misused. An audit can
be waved as empirical validation (``the sources were audited, so the model
works''); the boundary denies it. It can be read as physical identification
(``$V$ and $R$ were audited, so they measure something physical''); denied.
It can be marketed (``audited'' as a product adjective); denied. And it can
be treated as operational authorization (``the audit approves intervening on
a model''); denied. The record also carries its status fields---admission
ring 4, license status \textsc{not licensed}, disclosure status
\textsc{held}---so that the boundary travels with the bytes.

The status fields merit unpacking, since they recur on every artifact in the
candidate's evidence chain and their meanings are load-bearing. Admission
ring 4 places the record at the outermost ring of the project's admission
scheme: material that may be worked on and cited internally but that has
earned no promotion toward canonical status. \textsc{Not licensed} states
that no use of the record grounds a product, marketing, or capability
claim---the record can be true and complete and still license nothing,
because licensing is a separate act by a separate authority that has not
occurred. \textsc{Held} states that the record is not disclosed: its
existence and content are internal, and nothing in this monograph's
reproduction of it changes its disclosure posture, since the monograph
inherits the same status labels (see the front matter). Finally, the
record's document status---\textsc{draft / open / not frozen}, as of
2026-08-04---means the record itself may still be amended by its owners; this
chapter documents the record as of that date, and a reader holding a later
revision should trust the artifact's own version fields over this text.

\begin{boundary}\label{ch04:bound:record}
Everything this chapter asserts is bounded by the record's own claim
boundary. The chapter documents that eight conflicts were found, how they
were evidenced, and what typed rules the candidate adopts in response. It
establishes nothing about whether any CAS-Ohm prediction is empirically
accurate, whether any audited quantity is physical, or whether any use of the
framework is licensed. All of that remains \textsc{open / ring 4 / not
licensed / held}.
\end{boundary}

\section{Why Preserving Conflicts Is an Engineering Result}\label{ch04:preservation}

It may seem deflationary to end a chapter of conflicts with no resolutions.
The claim of this section is the opposite: for a framework in CAS-Ohm's
epistemic position, the preserved audit is worth more than any resolution
currently available, for five reasons.

\emph{First, harmonization without evidence is fabrication.} Every conflict
above could have been ``resolved'' editorially: declare $R$ to be $0$--$100$,
declare $I_{\mathrm{FS}}$ to be the corpus percentile, declare $\beta$ to be
the Jacobian derivative, declare the network rules valid. Each declaration
would have manufactured exactly the coherence the sources demonstrably lack,
and every downstream artifact would then cite ``the framework'' for choices
that were in fact one editor's unevidenced preference. The audit's refusal to
harmonize is a refusal to launder editorial choice into inherited authority.
Where a choice genuinely must be made to proceed (Q4 bins, the split
classification axes), the record makes it as a \emph{named candidate rule}
attached to a conflict---visible, attributable, and reversible---rather than
as a silent rewrite of the sources.

\emph{Second, preserved conflicts are executable.} Because each disposition
is implemented as a type-level constraint---distinct quantity kinds with no
silent conversion, an \texttt{IFSMethod} enumeration inside the profile
digest, separate unit kinds for scores and resistances, separately named gain
profiles, paired classification axes, capability gates on network arithmetic,
censoring rules for residuals, exact half-open bins---the audit is not a
document that future engineers must remember to consult. It is enforced at
validation time. The historically observed failure modes (the mixed-scale
$\mathrm{sat}=8.5$, the score-as-resistance division, the
residual-as-bypass leap) are not merely warned against; they are
unrepresentable or fail closed in the candidate's data model. An audit that
compiles is a different artifact from an audit that admonishes.

\emph{Third, the record converts ambiguity into falsifiable structure.}
Several conflicts are not decidable by inspection because they are empirical
questions in disguise: whether series and parallel rules hold, whether gains
compose, whether residuals localize faults. The audit routes each to a
registered hypothesis with a preregistered failure consequence (drop the
rule, withdraw the causal language). This is the correct fate for
conjectural machinery---neither adopted nor discarded, but staged for
controlled testing with the decision criteria fixed in advance.

\emph{Fourth, the audit is itself evidence about process.} Independent of
CAS-Ohm's eventual fortunes, the record demonstrates a reusable method for
inheriting a multi-document specification honestly: fix artifact identities
by digest; disclose locator systems and inspection procedures, including
failed ones; record disagreements verbatim with exact positions; classify the
failure kinds; attach dispositions as typed, named rules; and bind the whole
under an explicit claim boundary. Any successor framework, and any auditor of
this one, can check the work. The alternative process---silent
harmonization---leaves no residue to check.

\emph{Fifth, preservation is what mature measurement practice does with
unresolved definition.} Metrological convention distinguishes sharply
between a quantity, the procedure that realizes it, and the uncertainty of
its values \cite{vim,gum}; when two procedures realize different quantities
under one name, the metrological response is not to average them but to
name them separately and record the difference until a comparison experiment
relates them. The audit's dispositions---separate named profiles for
$I_{\mathrm{FS}}$ methods, separately named gain profiles, distinct quantity
types for scales---are this convention applied to an inherited specification
corpus. Nothing about the practice is novel; what would have been novel, and
wrong, is exempting a machine-learning-adjacent framework from it on the
grounds that its quantities are ``only scores.'' The audit treats score
scales with the same definitional seriousness a standards laboratory would
apply to a physical unit, precisely because the candidate's entire claim to
discipline rests on refusing the score/measurement blur documented in the
sources themselves.

The chapters that follow build directly on the dispositions established
here. Chapter~\ref{ch05:profile} develops the calibration profile whose
fields exist largely because Conflicts 1, 2, and 3 demand them;
Chapter~\ref{ch06:response} develops the response model and the
anti-circularity rule that Conflict 3's calibrated bridge makes possible and
that Conflict 7's discipline makes necessary. Throughout, the eight records
of this chapter function as the monograph's negative space: the places where
the sources do not speak with one voice, held open on purpose, so that
nothing built above them can quietly pretend that they do.
